\section*{Chapitre \ref{ch:ml:reinforcement-learning-for-execution-and-market-making} --- Apprentissage par renforcement pour l'exécution et le market making}

\begin{solution}{exo:ml:reinforcement-learning-for-execution-and-market-making:1}
Impact~: $10\times0.001\times0.1^2/0.1 = 0.001$. Risque~: $10\times0.02^2\times0.1\times(0.9^2 + 0.8^2 + \dots + 0.1^2) =
0.0004\times2.85 = 0.00114$. Total $0.00214$, soit 21.4 points de base.
\end{solution}

\begin{solution}{exo:ml:reinforcement-learning-for-execution-and-market-making:2}
$\frac{d}{d\delta}\delta e^{-k\delta} = (1 - k\delta)e^{-k\delta} = 0$ en $\delta = 1/k = 0.67$. La profondeur plate de l'optimum exact
au départ vaut 0.68~: les pénalités d'inventaire l'élargissent légèrement, puisque chaque exécution ajoute du risque d'inventaire.
\end{solution}

\begin{solution}{exo:ml:reinforcement-learning-for-execution-and-market-making:3}
La cible évalue l'état suivant par sa meilleure action~; si des actions interdites entrent dans le maximum, l'agent apprend des valeurs
qu'il ne pourra jamais réaliser (vendre plus qu'il ne détient, sauter la vente finale forcée), et ces valeurs gonflées se propagent
en remontant vers chaque état antérieur.
\end{solution}

\begin{solution}{exo:ml:reinforcement-learning-for-execution-and-market-making:4}
$Z$ est indépendant de tout ce que l'agent connaît et choisit, donc $\E[\sigma\sqrt\tau Zx_k\mid\text{historique}] = 0$ pour toute
\omterm{def:ml:reinforcement-learning-foundations:mdp}{politique}~: le rendement espéré de chaque \omterm{def:ml:reinforcement-learning-foundations:mdp}{politique} est inchangé, donc celui de la \omterm{def:ml:reinforcement-learning-foundations:mdp}{politique} optimale aussi. L'écart-type du terme, environ
$63x$ points de base par pas, écrase les écarts entre calendriers (quelques points de base en tout), de sorte que, avec le bruit, les
estimations de valeur du réseau ont été dominées par ce terme après 1\,500 épisodes, et son choix glouton était le calendrier plat~; avec plus
d'épisodes, il s'améliorerait, mais lentement.
\end{solution}

\begin{solution}{exo:ml:reinforcement-learning-for-execution-and-market-making:5}
Sa troisième entrée, l'impact observé, valait toujours zéro à l'entraînement (l'impact ne variait jamais), si bien que la réponse du réseau à
d'autres valeurs n'était jamais contrainte~; pour $\log3$, il a produit un calendrier arbitraire. La forme fermée du mauvais modèle est au
moins un calendrier raisonnable pour un monde voisin, et son erreur croît régulièrement avec l'erreur sur le paramètre.
\end{solution}

\begin{solution}{exo:ml:reinforcement-learning-for-execution-and-market-making:6}
La robustesse à la volatilité ne dit rien de ce qui n'a pas été randomisé (modèle d'exécution, sélection adverse, latence, autres teneurs de
marché) et rien sur la structure du simulateur~; 10\,000 simulations d'un même simulateur ne constituent qu'un seul élément de preuve. Tester dans des mondes
qui diffèrent sur les hypothèses non randomisées, et sur des journaux réels.
\end{solution}

\begin{solution}{exo:ml:reinforcement-learning-for-execution-and-market-making:7}
La recherche sur grille retient $c = 0.70$ et $d = 0.03$ sur 2\,000 trajectoires~; sur les 5\,000 trajectoires communes du chapitre, la \omterm{def:ml:reinforcement-learning-foundations:mdp}{politique}
obtient 67.82, contre 67.95 pour l'optimum exact et 62.84 pour la table de valeurs après 8\,000 épisodes.
Deux paramètres bien choisis retrouvent presque tout, ce qui plaide pour la structure plutôt que pour les tables.
\end{solution}

\begin{solution}{exo:ml:reinforcement-learning-for-execution-and-market-making:8}
Minimiser $\sum_{k=1}^n\eta(x_{k-1} - x_k)^2/\tau + \lambda\sigma^2\tau\sum_{k=1}^{n-1}x_k^2$ avec $x_0 = X$, $x_n = 0$. La
dérivée en $x_k$ donne $-2\eta(x_{k-1} - x_k)/\tau + 2\eta(x_k - x_{k+1})/\tau + 2\lambda\sigma^2\tau x_k = 0$, soit
$x_{k-1} - 2x_k + x_{k+1} = (\lambda\sigma^2\tau^2/\eta)x_k$. Les solutions sont des combinaisons de $e^{\pm\tilde\kappa t_k}$ avec
$2(\cosh\tilde\kappa\tau - 1) = \lambda\sigma^2\tau^2/\eta$~; les conditions aux limites donnent $x_k = X\sinh(\tilde\kappa(T -
t_k))/\sinh(\tilde\kappa T)$, le calendrier de \texttt{firm.acexec.discrete}.
\end{solution}

\begin{solution}{pb:ml:reinforcement-learning-for-execution-and-market-making:1}
\textbf{Partie I.}
\begin{enumerate}
\item État~: temps, positions détenues, impact observé~; actions~: lots à vendre (actions masquées)~; \omterm{def:ml:reinforcement-learning-foundations:mdp}{récompense}~: opposé du coût d'impact moins la charge de risque.
\item Son espérance est nulle, donc il n'ajoute que de la variance~; avec lui, le DQN a appris le TWAP (à 2.55 points de base de l'optimum).
\item Vendre plus que ce qui est détenu, et toute autre quantité que la totalité du reliquat à la dernière étape, aussi bien dans le choix que dans la cible.
\item C'est ce qu'un trader utiliserait pour s'adapter, et cela permet à une \omterm{def:ml:reinforcement-learning-foundations:mdp}{politique} entraînée sur de nombreux impacts d'agir sur celui qu'elle rencontre.
\end{enumerate}
\textbf{Partie II.}
\begin{enumerate}[resume]
\item TWAP 21.40, optimum sur grille de lots 19.05, DQN 19.29 points de base (optimum continu 18.85).
\item Monde du modèle, impact $\times3$, impact $/3$~: forme fermée du modèle 0, 2.20, 1.21~; TWAP 2.55, 1.04, 4.99~; DQN 0.44, 9.09,
4.60~; DQN avec \omterm{def:ml:reinforcement-learning-foundations:mdp}{récompense} bruitée 2.55, 1.04, 4.99~; DQN randomisé 0.55, 0.86, 1.24.
\item Avant sa première vente, il ne connaît pas l'impact, si bien que sa première étape se couvre sur toute la plage d'entraînement.
\item La fonction objectif gagnerait un terme de dérive~; un signal dans l'état permettrait à l'agent d'apprendre à s'en servir, et à
le surapprendre.
\end{enumerate}
\textbf{Partie III.}
\begin{enumerate}[resume]
\item 0.68 des deux côtés sans position~; 0.57 au cours vendeur et 0.79 au cours acheteur avec cinq unités en position longue.
\item Optimum exact 67.95~; Avellaneda--Stoikov 64.73~; $1/k$ constant 66.67~; \omterm{def:ml:reinforcement-learning-foundations:td}{Q-learning} 62.34 après 2\,000 épisodes et 62.84
après 8\,000.
\item Les espérances des profondeurs voisines diffèrent de moins que le bruit des exécutions qu'il observe~; chacune des 15\,120 entrées
de la table est estimée à partir de peu d'échantillons bruités.
\item Conserver la structure (profondeurs comme fonctions de l'inventaire et du temps, issues de la forme fermée ou d'une famille paramétrique) et apprendre
ce que le modèle omet~: sélection adverse, position dans la file, modèle d'exécution.
\end{enumerate}
\textbf{Partie IV.}
\begin{enumerate}[resume]
\item \emph{Battre une forme fermée.} Le DQN entraîné sur l'impact du modèle se situe à 0.44 point de base
de l'optimum d'Almgren--Chriss chez lui, et à 9.09 et 4.60 dans des mondes où l'impact est triplé ou divisé par trois~; entraîné avec \omterm{def:ml:reinforcement-learning-for-execution-and-market-making:dr}{randomisation de domaine},
et l'impact observé dans son état~: 0.55, 0.86 et 1.24.
\item Que ces méthodes fonctionnent dans des simulateurs, souvent ceux des auteurs, et dans quelques backtests sur données réelles~; pas la manière dont les \omterm{def:ml:reinforcement-learning-foundations:mdp}{politiques}
se sont comportées en conditions réelles, ce qui est rarement public.
\item Dans des limites de participation, avec la forme fermée en repli, comparée à la \omterm{def:ml:reinforcement-learning-foundations:mdp}{politique} en vigueur par \omterm{def:ml:reinforcement-learning-foundations:ope}{évaluation hors politique} (off-policy evaluation) et
un petit essai randomisé, et en surveillant son entrée d'impact observé.
\item Revenir à la forme fermée, avec l'impact réestimé~; ne pas laisser le réseau agir hors de la plage sur laquelle il a été entraîné.
\item L'état (positions dans la file, niveaux du carnet d'ordres, prix) est continu et de grande taille~; les tables ne généralisent pas.
\item Le \omterm{def:ml:reinforcement-learning-for-execution-and-market-making:dqn}{réseau cible} et le \omterm{def:ml:reinforcement-learning-for-execution-and-market-making:design}{masquage d'actions}~; retirer chacun tour à tour et comparer les courbes d'apprentissage selon les graines.
\item Comparaison appariée sur les mêmes ordres ou sur des ordres alternés, coûts mesurés par rapport au prix d'arrivée, avec assez d'ordres
pour que l'erreur type de la différence soit inférieure à l'effet.
\item Lorsque le problème comporte des parties qu'aucune forme fermée ne capture, et qu'un simulateur ou des données assez fidèles permettent de les apprendre.
\end{enumerate}
\end{solution}

\begin{solution}{iq:ml:reinforcement-learning-for-execution-and-market-making:1}
Le \omterm{def:ml:reinforcement-learning-for-execution-and-market-making:dqn}{rejeu d'expérience} brise la corrélation entre échantillons consécutifs et réutilise les données~; le \omterm{def:ml:reinforcement-learning-for-execution-and-market-making:dqn}{réseau cible} maintient la cible de régression fixe pendant
un certain temps, si bien que le réseau ne court pas après ses propres estimations mouvantes.
\iqlookfor{les deux mécanismes et l'instabilité que chacun traite.}
\end{solution}

\begin{solution}{iq:ml:reinforcement-learning-for-execution-and-market-making:2}
Les épisodes sont des ordres parents~; état~: temps restant, quantité restante, fourchette, profondeur, volatilité récente, position dans la file, impact observé~;
actions~: tailles et types des ordres enfants~; \omterm{def:ml:reinforcement-learning-foundations:mdp}{récompense}~: opposé des accroissements d'implementation shortfall, avec un terme de risque~; indices de référence~:
Almgren--Chriss et l'algorithme du desk.
\iqlookfor{l'état, l'action, la \omterm{def:ml:reinforcement-learning-foundations:mdp}{récompense} et un indice de référence.}
\end{solution}

\begin{solution}{iq:ml:reinforcement-learning-for-execution-and-market-making:3}
Modifier les \omterm{def:ml:reinforcement-learning-foundations:mdp}{récompenses} pour faciliter l'apprentissage sans changer la \omterm{def:ml:reinforcement-learning-foundations:mdp}{politique} optimale (retrait du bruit de moyenne nulle, termes fondés sur un potentiel). Cela tourne
mal lorsque la modification altère les incitations~: une prime sur les exécutions qui apprend à l'agent à franchir la fourchette, une pénalité d'inventaire qui
n'est pas le risque qui importe au desk.
\iqlookfor{la condition d'invariance et un exemple de façonnage nuisible.}
\end{solution}

\begin{solution}{iq:ml:reinforcement-learning-for-execution-and-market-making:4}
Surestimation due à la maximisation sur des estimations bruitées, éventuellement divergence due au bootstrap combiné à une approximation de fonction.
Utiliser le double \omterm{def:ml:reinforcement-learning-foundations:td}{Q-learning}, une mise à jour de la cible plus lente, un \omterm{def:ml:trees-and-boosting:boosting}{taux d'apprentissage} plus petit, une mise à l'échelle ou un écrêtage des \omterm{def:ml:reinforcement-learning-foundations:mdp}{récompenses}, et vérifier que les actions masquées
sont exclues de la cible.
\iqlookfor{le biais de maximisation et des remèdes concrets.}
\end{solution}

\begin{solution}{iq:ml:reinforcement-learning-for-execution-and-market-making:5}
Lorsque les formes fermées omettent des éléments importants (sélection adverse, dynamique de la file, plusieurs plateformes de négociation) et qu'il existe un bon simulateur ou de bons journaux,
avec la structure conservée dans la \omterm{def:ml:reinforcement-learning-foundations:mdp}{politique}. Non pour apprendre ce qu'une forme fermée donne déjà, où le signal par décision
est trop faible pour apprendre à partir des exécutions.
\iqlookfor{d'abord la structure, l'apprentissage pour les parties omises, et le problème du rapport signal sur bruit.}
\end{solution}

\begin{solution}{iq:ml:reinforcement-learning-for-execution-and-market-making:6}
Entraîner sur de nombreux tirages de paramètres du simulateur, de façon que la \omterm{def:ml:reinforcement-learning-foundations:mdp}{politique} fonctionne sur tous ces tirages. Cela aide lorsque les vrais paramètres se trouvent
dans la plage et que l'état permet à la \omterm{def:ml:reinforcement-learning-foundations:mdp}{politique} de les inférer~; cela nuit lorsque la plage est fausse ou trop large (la \omterm{def:ml:reinforcement-learning-foundations:mdp}{politique} se couvre et perd
chez soi) ou lorsque c'est la structure non randomisée qui est fausse.
\iqlookfor{ce qu'elle apporte, la condition d'observabilité, et ses coûts.}
\end{solution}
